\documentclass[12pt, a4paper, twoside]{report}

% ── Encoding & Language ───────────────────────────────────────────────────────
\usepackage[utf8]{inputenc}
\usepackage[T1]{fontenc}
\usepackage[italian]{babel}

% ── Typography ────────────────────────────────────────────────────────────────
\usepackage{lmodern}
\usepackage{microtype}
\usepackage{setspace}
\onehalfspacing

% ── Page Layout ───────────────────────────────────────────────────────────────
\usepackage[
  a4paper,
  top=3cm, bottom=3cm,
  inner=3.5cm, outer=2.5cm,
  headheight=15pt
]{geometry}

% ── Headers & Footers ─────────────────────────────────────────────────────────
\usepackage{fancyhdr}
\pagestyle{fancy}
\fancyhf{}
\fancyhead[LE]{\small\textit{Capitolo 1 — Teoria dei Computer Quantistici}}
\fancyhead[RO]{\small\textit{\nouppercase{\rightmark}}}
\fancyfoot[C]{\thepage}
\renewcommand{\headrulewidth}{0.4pt}

% ── Math ──────────────────────────────────────────────────────────────────────
\usepackage{amsmath, amssymb, amsthm}
\usepackage{braket}          % Dirac notation: \ket{}, \bra{}, \braket{}
\usepackage{physics}

% ── Tables ────────────────────────────────────────────────────────────────────
\usepackage{booktabs}
\usepackage{tabularx}
\usepackage{array}
\usepackage{xcolor}
\usepackage{colortbl}

% ── Graphics & Figures ────────────────────────────────────────────────────────
\usepackage{graphicx}
\usepackage{tikz}
\usepackage{quantikz}

% ── Listings (codice) ─────────────────────────────────────────────────────────
\usepackage{listings}
\usepackage{lstautogobble}

% ── Hyperlinks ────────────────────────────────────────────────────────────────
\usepackage[
  colorlinks=true,
  linkcolor=blue!60!black,
  citecolor=green!50!black,
  urlcolor=blue!60!black
]{hyperref}

% ── Custom Colors ─────────────────────────────────────────────────────────────
\definecolor{accent}{RGB}{31, 78, 140}
\definecolor{lightblue}{RGB}{214, 228, 240}
\definecolor{lightgray}{RGB}{245, 245, 245}
\definecolor{darkgray}{RGB}{68, 68, 68}

% ── Chapter Style ─────────────────────────────────────────────────────────────
\usepackage{titlesec}
\titleformat{\chapter}[display]
  {\normalfont\Huge\bfseries\color{accent}}
  {\chaptertitlename\ \thechapter}{20pt}{\Huge}
\titlespacing*{\chapter}{0pt}{50pt}{40pt}

\titleformat{\section}
  {\normalfont\Large\bfseries\color{accent}}
  {\thesection}{1em}{}[\vspace{-6pt}\textcolor{accent}{\rule{\linewidth}{0.6pt}}]

\titleformat{\subsection}
  {\normalfont\large\bfseries\color{accent!80!black}}
  {\thesubsection}{1em}{}

% ── Theorem Environments ──────────────────────────────────────────────────────
\theoremstyle{definition}
\newtheorem{definizione}{Definizione}[chapter]
\newtheorem{principio}{Principio}[chapter]
\theoremstyle{remark}
\newtheorem{osservazione}{Osservazione}[chapter]

% ── Utility box ───────────────────────────────────────────────────────────────
\usepackage{tcolorbox}
\tcbuselibrary{skins, breakable}
\newtcolorbox{notabox}{
  colback=lightblue!40, colframe=accent,
  fonttitle=\bfseries, title=Nota,
  breakable, left=6pt, right=6pt
}

% ─────────────────────────────────────────────────────────────────────────────
\begin{document}

% ── Fake front matter for standalone compilation ──────────────────────────────
\pagenumbering{arabic}
\setcounter{chapter}{0}   % chapter 1 will be \chapter → numbered 1

% ═════════════════════════════════════════════════════════════════════════════
\chapter{Teoria dei Computer Quantistici}
\label{chap:quantum}
% ═════════════════════════════════════════════════════════════════════════════

\begin{center}
  \large\textit{Pro e Contro rispetto ai Computer Tradizionali}
\end{center}
\vspace{1cm}

% ─────────────────────────────────────────────────────────────────────────────
\section{Introduzione al Paradigma Quantistico}
\label{sec:intro}
% ─────────────────────────────────────────────────────────────────────────────

Negli ultimi anni, il progresso nell'elaborazione delle informazioni ha
seguito la celebre \textit{Legge di Moore}, che descrive il raddoppio
approssimativo del numero di transistor per chip ogni due anni. Tuttavia,
man mano che le dimensioni dei componenti elettronici si avvicinano alla scala
atomica, i limiti fisici della computazione classica diventano sempre più
evidenti. In questo contesto emerge il paradigma della
\textbf{computazione quantistica}, che sfrutta i principi della meccanica
quantistica --- sovrapposizione, entanglement e interferenza --- per elaborare
le informazioni in modo radicalmente diverso rispetto ai computer tradizionali.

Un computer classico elabora le informazioni attraverso \textit{bit} binari, che
possono assumere i valori $0$ o $1$. Al contrario, un computer quantistico
utilizza i \textbf{qubit} (\textit{quantum bit}), che grazie al principio di
sovrapposizione possono esistere contemporaneamente in una combinazione di
$\ket{0}$ e $\ket{1}$ fino al momento della misurazione. Questa proprietà,
combinata con l'entanglement quantistico e le operazioni di interferenza,
conferisce ai computer quantistici un potenziale computazionale enormemente
superiore per determinate classi di problemi.

% ─────────────────────────────────────────────────────────────────────────────
\section{Fondamenti Teorici della Computazione Quantistica}
\label{sec:fondamenti}
% ─────────────────────────────────────────────────────────────────────────────

\subsection{Il Qubit e la Sovrapposizione Quantistica}
\label{subsec:qubit}

\begin{definizione}[Qubit]
  Un \textbf{qubit} è il sistema fisico a due livelli che costituisce l'unità
  fondamentale di informazione nella computazione quantistica. Il suo stato puro
  è descritto da un vettore nello spazio di Hilbert $\mathcal{H} \cong \mathbb{C}^2$:
  \begin{equation}
    \ket{\psi} = \alpha\ket{0} + \beta\ket{1}, \qquad
    \alpha, \beta \in \mathbb{C}, \quad |\alpha|^2 + |\beta|^2 = 1.
    \label{eq:qubit}
  \end{equation}
\end{definizione}

I coefficienti $\alpha$ e $\beta$ sono le \textit{ampiezze di probabilità}: una
misurazione restituisce $0$ con probabilità $|\alpha|^2$ e $1$ con probabilità
$|\beta|^2$, collassando lo stato in uno dei due vettori base. Geometricamente,
ogni stato puro di un qubit corrisponde a un punto sulla superficie della
\textbf{sfera di Bloch} (Figura~\ref{fig:bloch}), parametrizzata da:
\begin{equation}
  \ket{\psi} = \cos\frac{\theta}{2}\ket{0} + e^{i\varphi}\sin\frac{\theta}{2}\ket{1},
  \qquad \theta \in [0,\pi],\ \varphi \in [0, 2\pi).
\end{equation}

\begin{figure}[ht]
  \centering
  \begin{tikzpicture}[scale=1.6]
    % sphere outline
    \draw[thick, color=accent] (0,0) circle (1);
    \draw[dashed, color=accent!50] (-1,0) arc (180:360:1 and 0.3);
    \draw[color=accent!50] (-1,0) arc (180:0:1 and 0.3);
    % axes
    \draw[->, thick] (0,0) -- (0,1.35) node[above] {$\ket{0}$};
    \draw[->, thick] (0,0) -- (0,-1.35) node[below] {$\ket{1}$};
    \draw[->, thick] (0,0) -- (1.35,0) node[right] {$y$};
    \draw[->, thick] (0,0) -- (-0.7,-0.5) node[below left] {$x$};
    % state vector
    \draw[->, very thick, color=red!70!black] (0,0) -- (0.55,0.72)
      node[right] {$\ket{\psi}$};
    % angles
    \draw[dashed] (0,0.72) -- (0.55,0.72);
    \draw[dashed] (0,0) -- (0.55,0.72);
    \node at (0.15, 0.45) {\small $\theta$};
    \node at (0.35,-0.1) {\small $\varphi$};
  \end{tikzpicture}
  \caption{Rappresentazione geometrica di un qubit sulla sfera di Bloch. I poli
    nord e sud corrispondono rispettivamente agli stati base $\ket{0}$ e $\ket{1}$.}
  \label{fig:bloch}
\end{figure}

\subsection{L'Entanglement Quantistico}
\label{subsec:entanglement}

\begin{principio}[Entanglement]
  Due qubit si dicono \textbf{entangled} (intrecciati) se il loro stato
  composto $\ket{\Psi} \in \mathcal{H}^{\otimes 2}$ non è fattorizzabile come
  prodotto tensoriale degli stati individuali:
  $\ket{\Psi} \neq \ket{\psi_1} \otimes \ket{\psi_2}$.
\end{principio}

Il prototipo di stato entangled è fornito dalla base degli \textbf{stati di Bell}:
\begin{align}
  \ket{\Phi^+} &= \frac{1}{\sqrt{2}}\bigl(\ket{00} + \ket{11}\bigr), &
  \ket{\Phi^-} &= \frac{1}{\sqrt{2}}\bigl(\ket{00} - \ket{11}\bigr), \\
  \ket{\Psi^+} &= \frac{1}{\sqrt{2}}\bigl(\ket{01} + \ket{10}\bigr), &
  \ket{\Psi^-} &= \frac{1}{\sqrt{2}}\bigl(\ket{01} - \ket{10}\bigr).
\end{align}

L'entanglement è la risorsa centrale di algoritmi e protocolli quantistici
fondamentali quali la teletrasportazione quantistica, la distribuzione
quantistica delle chiavi (\textit{QKD}) e il \textit{dense coding}.
Einstein definì il fenomeno con l'espressione
\textit{``spooky action at a distance''}, ritenendolo inizialmente
incompatibile con il realismo locale; le disuguaglianze di Bell (1964) e le
successive verifiche sperimentali hanno confutato tale ipotesi.

\subsection{Le Porte Logiche Quantistiche}
\label{subsec:porte}

La computazione quantistica si effettua tramite \textbf{porte logiche
quantistiche}, ovvero trasformazioni unitarie $U \in \mathrm{U}(2^n)$ applicate
a registri di $n$ qubit. A differenza delle porte classiche, ogni operazione
quantistica (eccetto la misurazione) è \textit{reversibile}.
Le porte più importanti su singolo qubit sono:

\begin{equation}
  H = \frac{1}{\sqrt{2}}\begin{pmatrix}1 & 1 \\ 1 & -1\end{pmatrix}, \quad
  X = \begin{pmatrix}0 & 1 \\ 1 & 0\end{pmatrix}, \quad
  T = \begin{pmatrix}1 & 0 \\ 0 & e^{i\pi/4}\end{pmatrix}.
\end{equation}

La porta a due qubit \textbf{CNOT} è definita da:
\begin{equation}
  \mathrm{CNOT} = \begin{pmatrix}1&0&0&0\\0&1&0&0\\0&0&0&1\\0&0&1&0\end{pmatrix}.
\end{equation}

L'insieme $\{H, T, \mathrm{CNOT}\}$ è \textit{universale}: ogni trasformazione
unitaria può essere approssimata con precisione arbitraria da una composizione
finita di queste porte (teorema di Solovay--Kitaev). Un esempio di circuito per
generare lo stato di Bell $\ket{\Phi^+}$ è mostrato di seguito:

\begin{center}
\begin{quantikz}
  \lstick{$\ket{0}$} & \gate{H} & \ctrl{1} & \rstick[2]{$\ket{\Phi^+}$} \\
  \lstick{$\ket{0}$} & \qw      & \targ{}  &
\end{quantikz}
\end{center}

% ─────────────────────────────────────────────────────────────────────────────
\section{Vantaggi dei Computer Quantistici}
\label{sec:vantaggi}
% ─────────────────────────────────────────────────────────────────────────────

Il vantaggio computazionale dei sistemi quantistici si manifesta in specifiche
classi di problemi dove gli algoritmi quantistici offrono un'accelerazione
\textit{esponenziale} o \textit{polinomiale} rispetto ai migliori algoritmi
classici noti.

\begin{itemize}
  \item \textbf{Fattorizzazione e crittografia.}
    L'\textbf{algoritmo di Shor} (1994) fattorizza un intero $N$ in tempo
    polinomiale $O((\log N)^3)$, riducendo da esponenziale a polinomiale il
    costo rispetto ai migliori algoritmi classici. Questo ha implicazioni dirette
    sulla sicurezza dei sistemi crittografici RSA e ElGamal, basati sulla
    difficoltà computazionale classica della fattorizzazione.

  \item \textbf{Ricerca non strutturata.}
    L'\textbf{algoritmo di Grover} (1996) individua un elemento in un database
    non strutturato di $N$ voci in $O(\sqrt{N})$ valutazioni dell'oracolo,
    ottenendo uno \textit{speedup} quadratico rispetto alla ricerca classica
    $O(N)$. Sebbene meno spettacolare della fattorizzazione, questo vantaggio
    è applicabile a un'ampia classe di problemi.

  \item \textbf{Simulazione quantistica.}
    Feynman (1982) osservò che i sistemi quantistici sono naturalmente simulati
    da computer quantistici. La simulazione di molecole e reazioni chimiche con
    precisione quantomeccanica richiederebbe risorse classiche esponenziali in
    $n$ (numero di elettroni), ma è affrontabile efficientemente su hardware
    quantistico. Applicazioni dirette includono la progettazione di farmaci,
    catalizzatori e materiali per batterie.

  \item \textbf{Ottimizzazione combinatoria.}
    Algoritmi come \textit{QAOA} (Quantum Approximate Optimization Algorithm) e
    il \textit{Quantum Annealing} offrono approcci promettenti per problemi
    NP-difficili. Applicazioni in logistica, finanza e machine learning sono
    oggetto di ricerca attiva.

  \item \textbf{Crittografia quantistica.}
    La QKD (\textit{Quantum Key Distribution}), nella sua formulazione BB84
    (Bennett \& Brassard, 1984), garantisce la sicurezza informazione-teorica
    della distribuzione delle chiavi: qualsiasi tentativo di intercettazione
    altera irreversibilmente gli stati trasmessi, rendendosi rilevabile.
\end{itemize}

% ─────────────────────────────────────────────────────────────────────────────
\section{Limiti e Sfide dei Computer Quantistici}
\label{sec:limiti}
% ─────────────────────────────────────────────────────────────────────────────

Nonostante il loro potenziale rivoluzionario, i computer quantistici presentano
sfide tecniche e teoriche fondamentali che ne limitano attualmente
l'applicabilità pratica su larga scala.

\begin{itemize}
  \item \textbf{Decoerenza quantistica.}
    I qubit sono estremamente sensibili alle perturbazioni ambientali. La
    \textbf{decoerenza} --- la perdita della coerenza quantistica per interazione
    con l'ambiente --- distrugge le sovrapposizioni. I tempi di coerenza $T_1$
    (rilassamento) e $T_2$ (sfasamento) nei qubit superconduttori tipici sono
    dell'ordine di $10$--$500\,\mu$s, limitando la profondità dei circuiti
    eseguibili.

  \item \textbf{Correzione degli errori quantistici (QEC).}
    I \textit{codici di correzione degli errori quantistici} (codice di
    superficie, codice di Steane, \ldots) richiedono un elevato \textit{overhead}
    di qubit fisici per qubit logico. Stime recenti indicano
    $10^3$--$10^4$ qubit fisici per qubit logico affidabile. I computer attuali
    operano in regime \textbf{NISQ} (\textit{Noisy Intermediate-Scale Quantum}),
    privi di correzione degli errori completa.

  \item \textbf{Scalabilità hardware.}
    I processori superconduttori (IBM, Google) richiedono temperature criogeniche
    $\sim 15\,\text{mK}$, vicine allo zero assoluto. Le tecnologie alternative
    (qubit fotonici, trappole ioniche, qubit a spin in azoto vacante, qubit
    topologici) presentano ciascuna diversi compromessi tra tempo di coerenza,
    fedeltà delle porte e scalabilità.

  \item \textbf{Programmabilità e stack software.}
    Lo sviluppo di algoritmi quantistici richiede competenze in meccanica
    quantistica, complessità computazionale e architettura hardware. I framework
    attuali (Qiskit, Cirq, PennyLane, Braket) stanno maturando, ma la distanza
    rispetto all'accessibilità della programmazione classica rimane ampia.

  \item \textbf{Applicabilità circoscritta.}
    Per la grande maggioranza dei compiti computazionali quotidiani ---
    elaborazione di testi, database relazionali, navigazione web --- i computer
    classici rimangono più efficienti e pratici. Il vantaggio quantistico è
    limitato a specifiche classi di problemi strutturati.
\end{itemize}

\begin{notabox}
  Il concetto di \textbf{quantum advantage} (vantaggio quantistico) si riferisce
  alla dimostrazione che un dispositivo quantistico risolve un problema
  specifico più velocemente di qualsiasi computer classico noto. La sua
  dimostrazione sperimentale su problemi di rilevanza pratica rimane una delle
  sfide aperte più importanti del settore.
\end{notabox}

% ─────────────────────────────────────────────────────────────────────────────
\section{Confronto Sistemico: Classico vs.\ Quantistico}
\label{sec:confronto}
% ─────────────────────────────────────────────────────────────────────────────

La Tabella~\ref{tab:confronto} sintetizza le principali differenze tra le due
architetture computazionali.

\begin{table}[ht]
  \centering
  \renewcommand{\arraystretch}{1.35}
  \caption{Confronto tra computer classici e computer quantistici.}
  \label{tab:confronto}
  \begin{tabularx}{\textwidth}{>{\bfseries}lXX}
    \toprule
    \rowcolor{accent}
    \color{white}Caratteristica &
    \color{white}Computer Classico &
    \color{white}Computer Quantistico \\
    \midrule
    Unità base          & Bit ($0$ o $1$)               & Qubit ($\ket{0}$, $\ket{1}$, sovrapposizione) \\
    \rowcolor{lightgray}
    Operazioni          & Porte irreversibili            & Trasformazioni unitarie reversibili \\
    Parallelismo        & Sequenziale / multicore        & Parallelismo quantistico ($2^n$ ampiezze) \\
    \rowcolor{lightgray}
    Gestione errori     & Maturissima, affidabile        & In sviluppo (QEC, regime NISQ) \\
    Cond.\ operative    & Temperatura ambiente           & Criogenica ($\sim 15\,\text{mK}$) \\
    \rowcolor{lightgray}
    Applicaz.\ ottimali & Uso generale, I/O, database    & Crittografia, simulazione, ottimiz. \\
    Maturità tecnol.    & Estremamente matura            & Emergente \\
    \rowcolor{lightgray}
    Costi               & Accessibili, scala industriale & Elevati, infrastruttura specializzata \\
    \bottomrule
  \end{tabularx}
\end{table}

% ─────────────────────────────────────────────────────────────────────────────
\section{Stato dell'Arte e Prospettive Future}
\label{sec:stato}
% ─────────────────────────────────────────────────────────────────────────────

Nel 2019, Google ha annunciato di aver raggiunto la \textbf{supremazia
quantistica} con il processore \textit{Sycamore} da 53 qubit, completando in
$\sim 200\,\text{s}$ un campionamento casuale di circuiti che richiederebbe
stimativamente $10\,000$ anni al supercomputer \textit{Summit}. Tale affermazione
è stata contestata da IBM, che ha proposto algoritmi classici più efficienti per
lo stesso problema, evidenziando come la definizione di supremazia quantistica
richieda una precisazione del problema di riferimento.

Nel 2023, IBM ha presentato il processore \textbf{Condor} da $1\,121$ qubit,
mentre Microsoft ha accelerato lo sviluppo dei qubit topologici basati sui
fermioni di Majorana, che promettono maggiore resistenza intrínseca alla
decoerenza. La roadmap di IBM prevede processori con migliaia di qubit logici
entro il 2030.

Il modello di accesso dominante si sta spostando verso il paradigma
\textbf{cloud}: IBM Quantum Network, AWS Braket, Google Quantum AI e Azure
Quantum offrono accesso remoto a processori quantistici reali. Il percorso più
realistico per le applicazioni nei prossimi cinque-dieci anni è considerato il
paradigma \textbf{ibrido classico-quantistico}, in cui il computer quantistico
funge da co-processore acceleratore per specifici sottoproblemi, mentre il
controllo generale rimane classico.

% ─────────────────────────────────────────────────────────────────────────────
\section{Conclusioni del Capitolo}
\label{sec:conclusioni}
% ─────────────────────────────────────────────────────────────────────────────

La computazione quantistica rappresenta un cambio di paradigma fondamentale,
non una semplice evoluzione incrementale dei sistemi classici. I suoi vantaggi
teorici --- accelerazione esponenziale per la fattorizzazione, ricerca,
simulazione e ottimizzazione --- sono matematicamente dimostrati e
algoritmicamente fondati. Tuttavia, la distanza tra teoria e implementazione
hardware rimane considerevole: la decoerenza, gli errori fisici e le difficoltà
di scalabilità costituiscono ostacoli tecnici genuini e non ancora risolti.

Il futuro prevedibile non vede i computer quantistici \textit{sostituire} quelli
classici, bensì \textit{affiancarli} in architetture ibride specializzate.
I capitoli successivi analizzeranno in dettaglio le architetture hardware
esistenti, gli algoritmi quantistici più rilevanti e le prospettive applicative
nei settori della crittografia, della farmacologia computazionale e
dell'intelligenza artificiale.

\end{document}
